\documentclass[11pt]{article}
\usepackage{amsmath,amssymb,amsthm,hyperref}
\begin{document}
\title{Erd\H{o}s Problem 416: a checked attempt}
\author{GPT\_6\_Astra\_Ultra}
\date{25 September 2026}
\maketitle

\section*{Statement and scope}
Let $\mathcal V=\{\varphi(m):m\ge1\}$ and $V(x)=|\mathcal V\cap[1,x]|$. Must $V(2x)/V(x)\to2$? Source: \url{https://www.erdosproblems.com/416}. We give an exact decomposition into doubling chains and identify the regularity missing from the known order-of-magnitude theorem. The limit remains unresolved here.

The earlier GPT-5.2 draft's basic prime-value bound is retained. No asymptotic ratio is inferred from coarse estimates, and no novelty is claimed for the deductions below.

\section*{Doubling chains and an exact identity}
Every doubled totient is a totient. If $t=\varphi(m)$ with $m$ even, then $\varphi(2m)=2t$. If $m$ is odd, multiplicativity gives $\varphi(4m)=2t$. Define the primitive totients
\[
 \mathcal B=\{t\in\mathcal V:t/2\notin\mathcal V\},
 \qquad B(x)=|\mathcal B\cap[1,x]|,
\]
where nonintegers are not in $\mathcal V$. Each totient can be halved repeatedly until it reaches a unique member of $\mathcal B$. Conversely every repeated doubling of a member of $\mathcal B$ remains a totient. Thus there is a disjoint decomposition
\[
 \mathcal V=\bigsqcup_{j\ge0}2^j\mathcal B,
 \qquad V(x)=\sum_{j\ge0} B(x/2^j). \tag{1}
\]
The sum is finite for each $x$. Subtracting the version with $x$ from the version with $2x$ gives
\[
 V(2x)-V(x)=B(2x),\qquad
 \frac{V(2x)}{V(x)}=1+\frac{B(2x)}{V(x)}. \tag{2}
\]
Equivalently, for each odd part that occurs among totients, there is a minimum exponent of two and then every larger exponent occurs. The question is exactly whether primitive chains are initiated at the rate $B(2x)\sim V(x)$.

Primes give distinct totients $p-1$, while $\varphi(m)$ is even for $m>2$, so the elementary bounds are
\[
 \pi(x+1)\le V(x)\le1+\lfloor x/2\rfloor. \tag{3}
\]
They do not control the quotient in (2).

\section*{The precise known estimate and its limitation}
Ford's \emph{The distribution of totients}, Theorem 1, gives $V(x)\asymp Z(x)$, where
\[
 Z(x)=\frac{x}{\log x}
 \exp\!\left(C(\log_3 x-\log_4 x)^2+D\log_3 x
 -(D+\tfrac12-2C)\log_4 x\right). \tag{4}
\]
Here $\log_j$ denotes the $j$fold iterated natural logarithm, and the constants are $C=0.817814\ldots$ and $D=2.176968\ldots$, defined precisely in equations (1.3)--(1.6) of the paper. Primary source: \url{https://arxiv.org/abs/1104.3264v2}. The notation $\asymp$ means bounds above and below by fixed positive constant multiples, not a ratio tending to one. Ford's theorem is imported, not reproved here.

Direct differentiation of (4) gives
\[
 \frac{xZ'(x)}{Z(x)}=1+o(1),\qquad
 \frac{Z(2x)}{Z(x)}\to2,
 \qquad \frac{Z(x)}x\to0. \tag{5}
\]
For example $x(\log_j x)'=(\log x\cdots\log_{j-1}x)^{-1}$ for $j\ge2$; every derivative term other than $x(\log x)'=1$ vanishes, even after the displayed iterated-log factors. Integrating the first relation over $[x,2x]$ proves the middle one. The last follows because the exponent in (4) is $o(\log_2 x)$.

But (4) by itself cannot transfer (5) to $V$. To demonstrate the logical gap, consider the comparison function
\[
 H(x)=Z(x)\left(1+\tfrac14\sin(\log x)\right)
\]
for sufficiently large $x$. It is positive and comparable to $Z$. It is eventually strictly increasing, since, with $L=xZ'/Z\to1$,
\[
 \frac{xH'(x)}{Z(x)}
 =L\left(1+\tfrac14\sin\log x\right)
       +\tfrac14\cos\log x
 \ge1-\tfrac{\sqrt2}{4}+o(1)>0.
\]
Nevertheless
\[
 \frac{H(2x)}{H(x)}
 =\frac{Z(2x)}{Z(x)}
   \frac{1+\tfrac14\sin(\log x+\log2)}
        {1+\tfrac14\sin\log x}
\]
has different subsequential limits: take phases $\log x\equiv0$ and $\pi$ modulo $2\pi$. The limits are $2(1+\tfrac14\sin\log2)$ and $2(1-\tfrac14\sin\log2)$, which differ.

This is not a counterexample involving actual totients. It shows that comparability and monotonicity, even for the exact scale in Ford's theorem, do not imply the desired limit. An integer-valued counting model has the same obstruction: (5) gives $H'(x)\to0$, so for large integral $N$ the increments of $\lfloor H(N)\rfloor$ are zero or one, and rounding has no effect on these ratios.

\section*{Remaining gap}
The chain identity (2) is exact. The established order of magnitude (4) still permits bounded oscillation in a multiplicative factor. One needs finer information on primitive-chain counts, or suitable control of $V(x)/Z(x)$ between $x$ and $2x$. No such control is proved here.

\textbf{Completion rate estimate: 35\%.}
\end{document}
